% GENERATED FILE. Edit canonical lessons, then run npm run generate:books.
% canonical-source-hash: fnv1a64:cb15ffe7843d22db
% canonical-lessons: TE-C280-ruddu, TE-C280-anduko, TE-C280-ceppu, TE-C280-cercu, TE-C280-jama-ceyu

\chapter{To rub, To receive, To tell, To admit, To deposit}
\label{ch:te-to-rub-to-receive-to-tell-to-admit-to-deposit}
\hlchaptermodality{280}

\begin{chapteropening}
\textbf{By the end of this chapter:} \emph{I can say to rub, to receive, to tell, to admit and to deposit.}

Complete the last lesson of chapter 280: I can say to rub, to receive, to tell, to admit and to deposit.
\end{chapteropening}

\section[\texorpdfstring{\te{రుద్దు} (ruddu)}{రుద్దు (ruddu)}]{\te{రుద్దు} (ruddu) — to rub, to scrub}

\label{lesson:TE-C280-ruddu}



\begin{quote}
Before the new one: say the Telugu for to touch, then the Telugu for to peel.
\end{quote}

\subsection*{You'll want to know: \te{రుద్దు}}
\textbf{\te{రుద్దు}} — \emph{ruddu} — \textquotedblleft{}to rub, to scrub\textquotedblright{}.

\begin{grammarlens}[title={What you've built}]
One more word for this chapter.
\end{grammarlens}

\subsection*{Your turn}
\begin{itemize}
\raggedright
  \item \emph{Say it:} \emph{ruddu}
  \item \emph{Say it:} \emph{ruddu}, once more
  \item \emph{Say it:} say \emph{ruddu}
  \item \emph{Recall:} say the Telugu for to touch, then the Telugu for to peel, then say \emph{ruddu} again
\end{itemize}

\begin{tcolorbox}[breakable,colback=teal!4,colframe=teal!35!black,title={Before you move on}]
What does \te{రుద్దు} mean? (\textquotedblleft{}To rub, to scrub\textquotedblright{}.) Say it once more.
\end{tcolorbox}
\section[\texorpdfstring{\te{అందుకో} (andukō)}{అందుకో (andukō)}]{\te{అందుకో} (andukō) — to receive, to catch}

\label{lesson:TE-C280-anduko}



\begin{quote}
Before the new one: say the Telugu for to peel, then the Telugu for to rub.
\end{quote}

\subsection*{You'll want to know: \te{అందుకో}}
\textbf{\te{అందుకో}} — \emph{andukō} — \textquotedblleft{}to receive, to catch\textquotedblright{}.

\begin{grammarlens}[title={What you've built}]
One more word for this chapter.
\end{grammarlens}

\subsection*{Your turn}
\begin{itemize}
\raggedright
  \item \emph{Say it:} \emph{andukō}
  \item \emph{Say it:} \emph{andukō}, once more
  \item \emph{Say it:} say \emph{andukō}
  \item \emph{Recall:} say the Telugu for to peel, then the Telugu for to rub, then say \emph{andukō} again
\end{itemize}

\begin{tcolorbox}[breakable,colback=teal!4,colframe=teal!35!black,title={Before you move on}]
What does \te{అందుకో} mean? (\textquotedblleft{}To receive, to catch\textquotedblright{}.) Say it once more.
\end{tcolorbox}
\section[\texorpdfstring{\te{చెప్పు} (ceppu)}{చెప్పు (ceppu)}]{\te{చెప్పు} (ceppu) — to tell, to say}

\label{lesson:TE-C280-ceppu}



\begin{quote}
Before the new one: say the Telugu for to rub, then the Telugu for to receive.
\end{quote}

\subsection*{You'll want to know: \te{చెప్పు}}
\textbf{\te{చెప్పు}} — \emph{ceppu} — \textquotedblleft{}to tell, to say\textquotedblright{}.

\begin{grammarlens}[title={What you've built}]
One more word for this chapter.
\end{grammarlens}

\subsection*{Your turn}
\begin{itemize}
\raggedright
  \item \emph{Say it:} \emph{ceppu}
  \item \emph{Say it:} \emph{ceppu}, once more
  \item \emph{Say it:} say \emph{ceppu}
  \item \emph{Recall:} say the Telugu for to rub, then the Telugu for to receive, then say \emph{ceppu} again
\end{itemize}

\begin{tcolorbox}[breakable,colback=teal!4,colframe=teal!35!black,title={Before you move on}]
What does \te{చెప్పు} mean? (\textquotedblleft{}To tell, to say\textquotedblright{}.) Say it once more.
\end{tcolorbox}
\section[\texorpdfstring{\te{చేర్చు} (cērcu)}{చేర్చు (cērcu)}]{\te{చేర్చు} (cērcu) — to admit, to enrol, to add}

\label{lesson:TE-C280-cercu}



\begin{quote}
Before the new one: say the Telugu for to receive, then the Telugu for to tell.
\end{quote}

\subsection*{You'll want to know: \te{చేర్చు}}
\textbf{\te{చేర్చు}} — \emph{cērcu} — \textquotedblleft{}to admit, to enrol, to add\textquotedblright{}.

\begin{grammarlens}[title={What you've built}]
One more word for this chapter.
\end{grammarlens}

\subsection*{Your turn}
\begin{itemize}
\raggedright
  \item \emph{Say it:} \emph{cērcu}
  \item \emph{Say it:} \emph{cērcu}, once more
  \item \emph{Say it:} say \emph{cērcu}
  \item \emph{Recall:} say the Telugu for to receive, then the Telugu for to tell, then say \emph{cērcu} again
\end{itemize}

\begin{tcolorbox}[breakable,colback=teal!4,colframe=teal!35!black,title={Before you move on}]
What does \te{చేర్చు} mean? (\textquotedblleft{}To admit, to enrol, to add\textquotedblright{}.) Say it once more.
\end{tcolorbox}
\section[\texorpdfstring{\te{జమ} \te{చేయు} (jama cēyu)}{జమ చేయు (jama cēyu)}]{\te{జమ} \te{చేయు} (jama cēyu) — to deposit, to credit}

\label{lesson:TE-C280-jama-ceyu}



\begin{quote}
Before the new one: say the Telugu for to tell, then the Telugu for to admit.
\end{quote}

\subsection*{You'll want to know: \te{జమ} \te{చేయు}}
\textbf{\te{జమ} \te{చేయు}} — \emph{jama cēyu} — \textquotedblleft{}to deposit, to credit\textquotedblright{}.

\begin{grammarlens}[title={What you've built}]
One more word for this chapter.
\end{grammarlens}

\subsection*{Your turn}
\begin{itemize}
\raggedright
  \item \emph{Say it:} \emph{jama cēyu}
  \item \emph{Say it:} \emph{jama cēyu}, once more
  \item \emph{Say it:} say \emph{jama cēyu}
  \item \emph{Recall:} say the Telugu for to tell, then the Telugu for to admit, then say \emph{jama cēyu} again
\end{itemize}

\begin{tcolorbox}[breakable,colback=teal!4,colframe=teal!35!black,title={Before you move on}]
What does \te{జమ} \te{చేయు} mean? (\textquotedblleft{}To deposit, to credit\textquotedblright{}.) Say it once more.
\end{tcolorbox}
